% © 2026 Ing. David Jaroš — CC BY-NC-ND 4.0
%
% This work is licensed under a Creative Commons Attribution-NonCommercial-NoDerivatives
% 4.0 International License (CC BY-NC-ND 4.0).
%
% License History: Earlier drafts (up to v0.3) were released under CC BY 4.0.
% From v0.4 onward, all material is released under CC BY-NC-ND 4.0 to protect
% the integrity of the theoretical work during ongoing academic development.
%
% See LICENSE.md for full license text.

% bootstrap_step_m3_crossing.tex
% Modular Bootstrap Step M3: Crossing Symmetry Test
%
% Goal: Test whether s-channel = t-channel OPE for k=1, k=2, k→∞.
% Status: IN PROGRESS — see MODULAR_BOOTSTRAP_K1_PLAN.md §4 Step M3
% Prerequisite: bootstrap_step_m2_4point.tex (4-point function explicit)

\ifdefined\INCLUDEMODE
\else
  \documentclass[12pt,a4paper]{article}
  \usepackage[a4paper, margin=2.5cm]{geometry}
  \usepackage{amsmath, amssymb, amsthm}
  \usepackage{mathtools}
  \usepackage{hyperref}
  \usepackage{xcolor}
  \usepackage{mdframed}
  \usepackage{booktabs}
  \usepackage{tcolorbox}

  \newtheorem{theorem}{Theorem}[section]
  \newtheorem{lemma}[theorem]{Lemma}
  \newtheorem{proposition}[theorem]{Proposition}
  \newtheorem{corollary}[theorem]{Corollary}
  \theoremstyle{definition}
  \newtheorem{definition}[theorem]{Definition}
  \theoremstyle{remark}
  \newtheorem{remark}[theorem]{Remark}

  \newmdenv[backgroundcolor=yellow!20, linecolor=orange, linewidth=1.5pt]{gapbox}
  \newmdenv[backgroundcolor=green!15, linecolor=green!60!black, linewidth=1.5pt]{resultbox}
  \newmdenv[backgroundcolor=red!8,   linecolor=red!60,   linewidth=1.5pt]{warnbox}
  \newmdenv[backgroundcolor=blue!8,  linecolor=blue!50,  linewidth=1.5pt]{insightbox}

  \title{\textbf{Modular Bootstrap Step M3}\\
         \large Crossing Symmetry Test for $k = 1$, $k = 2$, $k \to \infty$}
  \author{Ing.\ David Jaro\v{s}}
  \date{2026-04-28}

  \begin{document}
  \maketitle
\fi

%──────────────────────────────────────────────────────────────────────────────
\section{Purpose and Status}
\label{sec:m3:purpose}
%──────────────────────────────────────────────────────────────────────────────

\begin{tcolorbox}[colback=blue!4!white, colframe=blue!50!black,
                  title=Step M3 — Crossing Symmetry Test]
\begin{tabular}{ll}
Parent plan  & \texttt{MODULAR\_BOOTSTRAP\_K1\_PLAN.md} \\
Prerequisite & \texttt{bootstrap\_step\_m2\_4point.tex}: 4-point function \\
Goal         & Determine which $k$ values satisfy crossing symmetry \\
Milestone    & 2026-05-19 \\
Status       & IN PROGRESS \\
Acceptance   & Unique $k$ forced, or explicit non-uniqueness \\
\end{tabular}
\end{tcolorbox}

\begin{warnbox}
\textbf{No-fit rule}: $k = 1$ must be \emph{forced} by crossing symmetry,
not assumed.  An argument showing $k = 1$ is ``simplest'' or ``natural''
has already been recorded as Motivated Conjecture [MC] (v67).  Step M3 must
go beyond [MC].
\end{warnbox}

%──────────────────────────────────────────────────────────────────────────────
\section{Crossing Symmetry Setup}
\label{sec:m3:crossing}
%──────────────────────────────────────────────────────────────────────────────

The 4-point function of Step M2 is:
\begin{equation}
  \mathcal{G}(z, \bar{z})
  \equiv \bigl\langle \Theta(\infty)\, \Theta(1)\, \Theta(z)\, \Theta(0)
         \bigr\rangle_{c=3}
  = |z|^{-3/2}\, |1 - z|^{-3/2}.
  \label{eq:m3:G4}
\end{equation}
Crossing symmetry requires:
\begin{equation}
  \mathcal{G}(z, \bar{z}) = \mathcal{G}(1-z, 1-\bar{z}).
  \label{eq:m3:crossing}
\end{equation}
Inspection of~\eqref{eq:m3:G4} shows this is satisfied:
$|z|^{-3/2}|1-z|^{-3/2} = |1-z|^{-3/2}|z|^{-3/2}$.
The full free-boson 4-point function is crossing symmetric for all $k$.
\emph{Crossing symmetry of the full theory does not fix $k$.}

\paragraph{The key question.}
Crossing symmetry of the 4-point function becomes a non-trivial constraint
when it is imposed on the \emph{conformal block decomposition} with a
\emph{restricted operator spectrum}.

\medskip
The three spectra to be tested:

\begin{center}
\renewcommand{\arraystretch}{1.3}
\begin{tabular}{clll}
\toprule
Label & Spectrum & Primaries & Condition \\
\midrule
$k = 1$ (SU(2)$_1$) & Two primaries per factor &
  $j = 0$ ($h=0$) and $j = 1/2$ ($h=1/4$) &
  No $h > 1/4$ primaries \\
$k = 2$ (SU(2)$_2$) & Three primaries per factor &
  $j = 0$ ($h=0$), $j=1/2$ ($h=3/8$), $j=1$ ($h=2/3$) &
  $j=1$ primary present \\
$k \to \infty$ (free boson) & Full winding tower &
  $h = (m+n)^2/4$, all $m,n \in \mathbb{Z}$ &
  Infinite tower \\
\bottomrule
\end{tabular}
\end{center}

%──────────────────────────────────────────────────────────────────────────────
\section{Conformal Block Decomposition with Restricted Spectrum}
\label{sec:m3:blocks}
%──────────────────────────────────────────────────────────────────────────────

The crossing symmetry equation in the conformal block basis is:
\begin{equation}
  \sum_{\mathcal{O} \in \mathcal{S}} C_{\mathcal{O}}^2\,
  \mathcal{F}_{\Delta_{\mathcal{O}}, \ell_{\mathcal{O}}}(z)
  \overline{\mathcal{F}}_{\Delta_{\mathcal{O}}, \ell_{\mathcal{O}}}(\bar{z})
  =
  \sum_{\mathcal{O} \in \mathcal{S}} C_{\mathcal{O}}^2\,
  \mathcal{F}_{\Delta_{\mathcal{O}}, \ell_{\mathcal{O}}}(1-z)
  \overline{\mathcal{F}}_{\Delta_{\mathcal{O}}, \ell_{\mathcal{O}}}(1-\bar{z}),
  \label{eq:m3:crossing_blocks}
\end{equation}
where the sum is over the spectrum $\mathcal{S}$ of allowed operators.

\subsection{Test 1: $k \to \infty$ (full free-boson spectrum)}

With the full winding tower, equation~\eqref{eq:m3:crossing_blocks} is
automatically satisfied (as shown in Step M2).
\begin{center}
\fbox{$k \to \infty$: crossing satisfied. \quad \textbf{CONSISTENT.}}
\end{center}

\subsection{Test 2: $k = 1$ (SU(2)$_1$ spectrum per factor)}

With only $j = 0, 1/2$ primaries (dimensions $h = 0$ and $h = 1/4$ per
Im($\mathbb{H}$) factor), the 4-point function truncates to:
\begin{equation}
  \mathcal{G}_{k=1}(z) = 1 + C_{1/2}^2\,\mathcal{F}_{1/2}(z)\overline{\mathcal{F}}_{1/2}(\bar{z})
  + \ldots \quad (\text{descendants included}).
  \label{eq:m3:G_k1}
\end{equation}
For $SU(2)_1$, the OPE coefficient is fixed by the WZW fusion rules:
\begin{equation}
  j_1 = j_2 = \tfrac{1}{2}, \quad j_3 = 0 \text{ or } j_3 = 1 \text{ (absent: } j \leq k/2 = 1/2 \text{)}.
  \label{eq:m3:fusion_k1}
\end{equation}
The fusion rule $[1/2] \times [1/2] = [0]$ (only identity survives the
$k=1$ truncation) gives $C_{j=1/2}^2 = 1$ and no $j=1$ contribution.

The crossing symmetry equation then becomes:
\begin{equation}
  \mathcal{F}_0(z) + C_{1/2}^2\,\mathcal{F}_{1/2}(z)
  =
  \mathcal{F}_0(1-z) + C_{1/2}^2\,\mathcal{F}_{1/2}(1-z).
  \label{eq:m3:crossing_k1}
\end{equation}
For $c = 1$ (one SU(2)$_1$ factor), the conformal blocks at $c = 1$ with
$h_{\mathrm{ext}} = 1/4$:
\begin{equation}
  \mathcal{F}_0(z) = z^{-1/2}(1 + O(z)), \quad
  \mathcal{F}_{1/2}(z) = z^0(1 + O(z)).
  \label{eq:m3:blocks_k1}
\end{equation}
At $z \to 0$: the $s$-channel is dominated by $\mathcal{F}_0 \sim z^{-1/2}$;
the $t$-channel ($z \to 1$) is dominated by $\mathcal{F}_0(1-z) \sim (1-z)^{-1/2}$.
The function $|z|^{-1/2}|1-z|^{-1/2}$ (the $c=1$ free-boson answer) is
reproduced exactly by the $k=1$ SU(2)$_1$ block sum, since the SU(2)$_1$ WZW
theory \emph{is} the free compact boson at the self-dual radius.
\begin{center}
\fbox{$k = 1$: crossing satisfied. \quad \textbf{CONSISTENT.}}
\end{center}

\subsection{Test 3: $k = 2$ (SU(2)$_2$ spectrum per factor)}

With $k = 2$, an extra primary $j = 1$ with $h = j(j+1)/(k+2) = 2/4 = 1/2$
(dimension $\Delta = 1$) is present.  The fusion rule $[1/2] \times [1/2]$
now includes both $[0]$ and $[1]$, giving an extra OPE contribution.
The crossing symmetry equation gains a new term:
\begin{equation}
  \mathcal{F}_0 + C_{1/2}^2 \mathcal{F}_{1/2} + C_1^2 \mathcal{F}_1
  =
  \mathcal{F}_0(1-z) + C_{1/2}^2 \mathcal{F}_{1/2}(1-z) + C_1^2 \mathcal{F}_1(1-z).
  \label{eq:m3:crossing_k2}
\end{equation}
The $SU(2)_2$ WZW OPE coefficient $C_1^2$ is non-zero.  The 4-point function
is \emph{not} equal to $|z|^{-3/2}|1-z|^{-3/2}$ (the free-boson / $k \to \infty$
answer): it contains the extra $j=1$ block at $c=3/2$ ($c=3/2$ per factor,
$c=9/2$ total — this is inconsistent with $c=3$).

\begin{warnbox}
\textbf{$k=2$ inconsistency}: $SU(2)_2$ has central charge
$c = 3k/(k+2) = 3 \times 2/4 = 3/2$ per factor, giving $c_{\mathrm{total}} = 9/2
\neq 3$.  The UBT partition function (Step M1) has $c = 3$.  Therefore $k = 2$
(and all $k \geq 2$) is \textbf{inconsistent} with the UBT central charge.
\end{warnbox}

\medskip
More generally, $SU(2)_k$ has central charge $c = 3k/(k+2)$.  Setting $c = 1$
per $\operatorname{Im}(\mathbb{H})$ factor:
\begin{equation}
  \frac{3k}{k+2} = 1 \quad \Longleftrightarrow \quad 3k = k + 2
  \quad \Longleftrightarrow \quad k = 1.
  \label{eq:m3:c_constraint}
\end{equation}

\begin{resultbox}
\textbf{Central-charge constraint}: The requirement $c = 1$ per
$\operatorname{Im}(\mathbb{H})$ factor (equivalently $c_{\mathrm{total}} = 3$,
proved in Step M1) forces $k = 1$ for the $SU(2)$ interpretation.
For $k \geq 2$: $c = 3k/(k+2) > 1$, inconsistent with $c_{\mathrm{total}} = 3$.
For $k \to \infty$ (free boson): $c \to 3$ per factor — but this is the $c=3$
free-compact-boson theory, which is equivalent to three $U(1)$ WZW models at
level $k \to \infty$, not $SU(2)_k$.
\end{resultbox}

%──────────────────────────────────────────────────────────────────────────────
\section{Crossing Summary and Gap Status}
\label{sec:m3:summary}
%──────────────────────────────────────────────────────────────────────────────

\begin{center}
\renewcommand{\arraystretch}{1.3}
\begin{tabular}{clll}
\toprule
$k$ & Crossing satisfied? & Consistent with $c=3$? & Verdict \\
\midrule
1 (SU(2)$_1$) & Yes & Yes ($c=1$ per factor) & \textbf{CONSISTENT} \\
2 (SU(2)$_2$) & Formally yes & No ($c=3/2$ per factor $\neq 1$) & \textbf{EXCLUDED} \\
$k \geq 2$     & Formally yes & No ($c = 3k/(k+2) > 1$) & \textbf{EXCLUDED} \\
$k \to \infty$ (free boson) & Yes & Yes ($c=1$ per factor) & \textbf{CONSISTENT} \\
\bottomrule
\end{tabular}
\end{center}

\paragraph{Residual degeneracy.}
Both $k = 1$ (SU(2)$_1$) and $k \to \infty$ (free compact boson) are consistent
with $c = 3$.  At the self-dual radius, $SU(2)_1$ WZW \emph{is} the free compact
boson at $R = \sqrt{\alpha'}$; they are the same CFT.  The degeneracy between
$k=1$ and $k \to \infty$ is resolved by the \emph{UBT operator content}, which
restricts to a finite set of primaries (Step M4).

\begin{gapbox}
\textbf{Remaining gap for Step M4}: The crossing symmetry + central-charge
argument excludes $k \geq 2$ but cannot distinguish $k = 1$ SU(2)$_1$
from $k \to \infty$ free boson, since these are equivalent at the self-dual
radius.  Step M4 must use the UBT field content (mode truncation from
N\_eff = 12, winding-number stability) to break this degeneracy.
\end{gapbox}

\ifdefined\INCLUDEMODE
\else
  \end{document}
\fi
